\section{LUT-Based Key Classification}
\label{sec:lut_classification}

\subsection{Motivation}

The final system does not solve a full continuous 3D inverse magnetic-field equation during runtime. Instead, it uses a measured lookup table (LUT). This approach is practical for our application because the piano keys form a discrete set of positions. The classifier only needs to determine which key is most likely, not estimate a continuous 3D coordinate.

Each key position produces a characteristic distribution of $H^2$ across the three sensors. The FPGA compares the live distribution with stored LUT entries and selects the key with the smallest error.

\subsection{LUT Construction}

For each frequency, a separate LUT is collected:

\begin{itemize}
    \item \code{h75\_lut.csv} for the 75 Hz finger.
    \item \code{h45\_lut.csv} for the 45 Hz finger.
\end{itemize}

The local sensing window contains five keys:

\begin{equation}
    k_{\mathrm{local}} \in \{-2,-1,0,+1,+2\}.
\end{equation}

For each key, three horizontal sample points are collected: top, middle, and bottom. To improve robustness to different finger heights, three vertical layers are collected: the key plane and two hover planes. Thus, one frequency LUT contains:

\begin{equation}
    5~\mathrm{keys} \times 3~\mathrm{points/key} \times 3~\mathrm{layers}
    = 45~\mathrm{entries}.
\end{equation}

Each entry is an average of multiple received $H^2$ frames, so the LUT represents a stable measured magnetic pattern.

\placeholderfigure{Insert the LUT sampling diagram: five local keys, three sample points per key, and three z layers.}{LUT sampling positions.}{fig:lut_sampling}{figures/lut_sampling.png}

\subsection{Normalized Feature Extraction}

Directly comparing raw $H^2$ values is sensitive to total field strength. If the finger is slightly closer or farther from the sensors, all $H^2$ values can change even when the position is the same. To focus on position, the classifier uses normalized sensor ratios:

\begin{align}
    H_{\mathrm{total}}^2 &= H_1^2 + H_2^2 + H_3^2, \\
    f_1 &= \frac{H_1^2}{H_{\mathrm{total}}^2}, \\
    f_2 &= \frac{H_2^2}{H_{\mathrm{total}}^2}, \\
    f_3 &= \frac{H_3^2}{H_{\mathrm{total}}^2}.
\end{align}

The ratios satisfy

\begin{equation}
    f_1 + f_2 + f_3 = 1.
\end{equation}

This means the classifier compares which sensor is strongest relative to the others, rather than comparing only absolute magnitude.

\subsection{Why Normalization Helps Position Classification}

Raw $H^2$ mainly describes signal strength. Normalized $H^2$ describes the spatial pattern across sensors. For example, when the magnetic finger is closer to sensor 1, sensor 1 becomes dominant. When it moves toward sensor 2, sensor 2 becomes more dominant. This relative distribution is more directly related to position.

A simple example is shown in Table~\ref{tab:normalization_example}. The live sample is actually Key A. A raw $H^2$ comparison would incorrectly choose Key B because the absolute values are closer to Key B. However, after normalization, the live pattern exactly matches Key A.

\begin{table}[H]
    \centering
    \caption{Example showing why normalized $H^2$ improves position classification.}
    \label{tab:normalization_example}
    \renewcommand{\arraystretch}{1.2}
    \begin{tabular}{l c c c c}
        \toprule
        \textbf{Sample} & $H_1^2$ & $H_2^2$ & $H_3^2$ & \textbf{Normalized pattern} \\
        \midrule
        LUT Key A & 10 & 5 & 5 & [0.50, 0.25, 0.25] \\
        LUT Key B & 18 & 9 & 3 & [0.60, 0.30, 0.10] \\
        Live at Key A & 20 & 10 & 10 & [0.50, 0.25, 0.25] \\
        \bottomrule
    \end{tabular}
\end{table}

The raw squared error to Key A is

\begin{equation}
    (20-10)^2 + (10-5)^2 + (10-5)^2 = 150.
\end{equation}

The raw squared error to Key B is

\begin{equation}
    (20-18)^2 + (10-9)^2 + (10-3)^2 = 54.
\end{equation}

Raw $H^2$ therefore chooses Key B incorrectly. In contrast, the normalized error to Key A is zero, so normalized features correctly preserve the position pattern.

\subsection{Error Function}

The FPGA uses the same core scoring method as the Python prototype. For a live feature vector and one LUT entry:

\begin{align}
    E_{\mathrm{ratio}} &=
    (f_{1,\mathrm{live}} - f_{1,\mathrm{LUT}})^2 +
    (f_{2,\mathrm{live}} - f_{2,\mathrm{LUT}})^2 \nonumber \\
    &\quad +
    (f_{3,\mathrm{live}} - f_{3,\mathrm{LUT}})^2, \\
    E_{\mathrm{strength}} &=
    \left(\ln(H_{\mathrm{total,live}}^2) -
    \ln(H_{\mathrm{total,LUT}}^2)\right)^2, \\
    E_{\mathrm{final}} &= E_{\mathrm{ratio}} + 0.05 E_{\mathrm{strength}}.
\end{align}

The logarithmic strength term keeps useful distance information without allowing raw magnitude to dominate the score.

\subsection{FPGA LUT ROM Format}

The LUT is stored in FPGA ROM initialized from a hexadecimal \code{.mem} file. Each line is one 160-bit LUT entry:

\begin{lstlisting}
[159:156] key_id
[155:128] reserved
[127:96]  f1_q30
[95:64]   f2_q30
[63:32]   f3_q30
[31:0]    ln_total_q16
\end{lstlisting}

The ratio values are stored in Q30 fixed-point format:

\begin{equation}
    f_{\mathrm{q30}} = f \cdot 2^{30}.
\end{equation}

The strength term is stored as Q16:

\begin{equation}
    \ln(H_{\mathrm{total}}^2)_{\mathrm{q16}}
    = \ln(H_{\mathrm{total}}^2) \cdot 2^{16}.
\end{equation}

During classification, the FPGA scans all 45 LUT entries. For each key, it keeps the minimum entry score among all sample points and z layers. The predicted local key is the key with the lowest final score.

\subsection{Classifier State Machine Implementation}

The classifier is implemented in \code{h2\_lut\_classifier\_3sensor.sv} as a multi-cycle finite-state machine. When \code{start} is asserted, the module latches the three live $H^2$ values and computes their total. It then uses the shared \code{unsigned\_divider.sv} block three times to calculate the Q30 ratios $f_1$, $f_2$, and $f_3$. The logarithmic strength feature is approximated in hardware by finding the base-2 exponent of the total strength and adding a small mantissa lookup correction in Q16.

After the live feature vector is ready, the classifier reads the 160-bit ROM entries initialized from the \code{.mem} LUT file. For each entry, it extracts the key ID, three Q30 ratio fields, and Q16 log-strength field. The state machine computes absolute differences, squares one component per cycle, adds the three ratio scores, adds the weighted strength score, and updates the best score stored for that key. After all 45 entries have been scanned, a final pass compares the five per-key best scores and outputs the minimum-score key with a one-clock \code{valid} pulse.
